\section{Computational Comparison Between Programming Paradigms}

The comparison between programming paradigms cannot be reduced solely to the
execution time of a program. A meaningful analysis must consider multiple
dimensions of software and computational behavior.

In this work, the comparison includes both quantitative and qualitative
criteria.

The quantitative metrics considered include:

\begin{itemize}
    \item execution time;
    \item memory usage;
    \item computational resource usage;
    \item number of iterations;
    \item implementation size;
    \item influence of compiler optimization.
\end{itemize}

The qualitative characteristics considered include:

\begin{itemize}
    \item readability;
    \item modularity;
    \item expressiveness;
    \item maintainability;
    \item correspondence between the mathematical formulation and its
    computational implementation.
\end{itemize}

The comparison must also account for the fact that C, Python, and Haskell
possess substantially different execution models.

C is generally compiled ahead of time into machine code and provides relatively
direct control over memory organization, data representation, and low-level
operations \cite{kernighan1988c}.

Python, in its most widely used implementation, CPython, compiles source code
into bytecode executed by the Python virtual machine. At the same time,
performance-critical scientific operations are frequently delegated to native
compiled libraries. This distinction is especially relevant in numerical
computing, where libraries such as NumPy provide optimized array operations
implemented largely outside ordinary Python execution
\cite{harris2020array}.

Haskell implementations such as the Glasgow Haskell Compiler typically compile
programs to native code while preserving semantic properties associated with
pure functional programming and non-strict evaluation. Consequently, the
runtime behavior of Haskell differs substantially from that of both C and
Python
\cite{marlow2010haskell}.

Experimental differences among implementations therefore cannot be attributed
exclusively to the programming paradigm. Observed performance results emerge
from the interaction among the numerical algorithm, programming paradigm,
programming language, compiler or interpreter, runtime system, optimization
strategies, memory representation, and hardware architecture.

For this reason, the experiments conducted in this research attempt to keep
external conditions constant whenever possible and explicitly document the
tools, language versions, compiler versions, compiler optimization options,
and hardware characteristics used during benchmarking.

Such methodological control is necessary to reduce confounding factors and to
avoid generalizing language-specific performance results into broader claims
about entire programming paradigms.

The objective of the experimental comparison is therefore not to determine an
absolute superiority between imperative and functional programming. Instead,
the purpose is to identify how their different computational models interact
with the mathematical and iterative properties of numerical algorithms.

This approach allows performance, memory behavior, implementation structure,
and numerical characteristics to be analyzed together, providing a broader
assessment of the consequences of programming-paradigm selection in scientific
and numerical computing.
